\section*{第 \ref{ch:b3:phase-transitions} 章——相变}

\begin{solution}{exo:b3:phase-transitions:1}
(a) 一级\omterm{def:b3:phase-transitions:order}{相变}：密度发生跳变，有潜热流动；
液体与气体之间没有对称性差别——这正是它们的\omterm{def:b3:hamiltonian-mechanics:portrait}{分界线}
可以\emph{终止}。(b) 连续\omterm{def:b3:phase-transitions:order}{相变}；\omterm{def:b3:phase-transitions:order}{序参量}为 $m$；上下（以及
转动）\omterm{def:b3:phase-transitions:order}{对称性破缺}。(c) 连续\omterm{def:b3:phase-transitions:order}{相变}；凝聚体的
宏观波函数（相角破缺）。(d) 连续\omterm{def:b3:phase-transitions:order}{相变}：
在\omterm{prop:b3:phase-transitions:vdw}{临界点}处跃变已缩为零——这是唯一一种
没有潜热的沸腾。
\end{solution}

\begin{solution}{exo:b3:phase-transitions:2}
(a) $Lm_{\text{w}}/Tk_{\text{B}} = 13\,k_{\text{B}}$，每个分子。
(b) $2.7\,k_{\text{B}}$。(c) 大致为 $k_{\text{B}}\ln$
（每个分子的千倍体积增量，$\ln1000 \approx 7$）加上
释放出的取向\omterm{def:b3:lagrangian-mechanics:coordinates}{自由度}：这正是特鲁顿规则中近乎普适的“十左右”。(d)
熔化仅略微释放取向与位置\omterm{def:b3:lagrangian-mechanics:coordinates}{自由度}——
分子仍紧密堆积；熵的大幅增加主要来自蒸气的
体积。
\end{solution}

\begin{solution}{exo:b3:phase-transitions:3}
(a) $3/8$，与物质种类无关——这是模型的对应态定律；
真实流体的 $0.29$ 表明范德瓦耳斯模型只是粗略的近似，
在\omterm{prop:b3:phase-transitions:vdw}{临界点}附近关联效应起主导作用。(b) $b =
k_{\text{B}}T_{\text{c}}/8P_{\text{c}} = \qty{7.1e-29}{m^3}$；$a =
27b^2P_{\text{c}} = \qty{1.0e-48}{J.m^3}$。(c) $b^{1/3} \approx
\qty{0.41}{nm}$——即分子的尺寸，由水蒸气
表求得。(d) 仅靠加压使气体液化需要 $T <
T_{\text{c}}$：在室温下对 CO$_2$ 成立，对
N$_2$ 则不成立——因此 CO$_2$ 钢瓶中装的是液体，氮气
钢瓶中只能装压缩气体。
\end{solution}

\begin{solution}{exo:b3:phase-transitions:4}
(a) $m$。(b) $\rho_{\text{液}} - \rho_{\text{气}}$。(c) 凝聚体的
振幅 $\psi_0$（连同其相位）。(d)
\emph{交错}磁化强度——两个
子晶格磁矩之差。
\end{solution}

\begin{solution}{exo:b3:phase-transitions:5}
(a) 单个自旋处于 $B_{\text{有效}} \propto qJm$ 中：$\langle s\rangle =
\tanh(\beta qJm)$；自洽性使循环闭合。(b) 非零
交点存在的充要条件是初始斜率 $\beta qJ > 1$。(c) $m =
\beta qJm - \tfrac13(\beta qJm)^3$ 得 $m^2 = 3(1 -
T/T_{\text{c}})$，在 $T_{\text{c}}$ 附近。(d) $J =
k_{\text{B}}T_{\text{c}}/q \approx \qty{11}{meV}$ 每个键——
这是电子伏特量级的交换作用，其来源见
\cref{exo:b3:identical-particles:12}，由八个近邻分担：
使铁的居里点升至一千开尔文的，是静电作用，
而非磁作用。
\end{solution}

\begin{solution}{exo:b3:phase-transitions:6}
(a) 线性化后，$m = \beta(qJm + \mu B)$：$\chi =
\mu\beta/(1 - T_{\text{c}}/T) \propto 1/(T - T_{\text{c}})$。(b)
两者都给出直线 $1/\chi$；自由自旋的直线
过原点，铁磁体的直线则在轴上截距为
$+T_{\text{c}}$：截距\emph{正是}耦合
强度的判据。(c) 镍在 \qty{631}{K} 时有序化：其高温
磁化率已预示了它将到达的居里点。
(d) 截距为负意味着 $J < 0$：近邻偏好
反平行排列——一种有待实现的反铁磁体。
\end{solution}

\begin{solution}{exo:b3:phase-transitions:7}
(a) $\dd P/\dd T = L/T\Delta v = \qty{3.6}{kPa/K}$。(b)
$\ln(0.7) = -4900(1/T - 1/373)$：$T \approx \qty{363}{K} =
\qty{90}{\celsius}$。(c) $\dd T/\dd P = T\Delta v/L =
\qty{-7.4e-8}{K/Pa}$：要降低五开尔文需要 $\sim\num{7e7}\,\unit{Pa}$
——近七百个大气压。(d) 滑冰者施加的
$\sim\qty{2e7}{Pa}$ 仅能带来 $-1.7$ K 的变化：加压熔化
在寒冷的冰场上行不通；润滑膜由摩擦生热
以及冰本身固有的预熔表层形成。
\end{solution}

\begin{solution}{exo:b3:phase-transitions:8}
(a) $m^2 = a(T_{\text{c}} - T)/2b$。(b) $F_{\min} =
-a^2(T_{\text{c}} - T)^2/4b$：$C = -T\partial^2F/\partial T^2$
在\omterm{def:b3:phase-transitions:order}{相变}处跃变 $a^2T_{\text{c}}/2b$。(c) \omterm{thm:b3:phase-transitions:meanfield}{平均场理论}：
一个干净的阶跃；氦的实测峰呈 $\lambda$ 形、近似对数型的
峰正是朗道所忽略的涨落所致——这一形状
给该\omterm{def:b3:phase-transitions:order}{相变}命名。(d) 加上 $-hm$ 后，对称性
已预先破缺：$m$ 永不恰好为零，每一个奇异性都
被抹平——自发破缺需要无偏的选择。
\end{solution}

\begin{solution}{exo:b3:phase-transitions:9}
(a) $\dd/\dd r$ 作用于 $4\pi r^2\sigma - \tfrac43\pi r^3n\delta\mu = 0$
在 $r^* = 2\sigma/n\delta\mu$ 处为零。(b) 小胚芽会收缩（表面
占优）：纯净相可以处于亚稳态等待，直至涨落
越过势垒；尘埃、离子和划痕提供了现成的表面，
可大幅降低势垒。(c) 离子径迹沿粒子路径催生液滴（云
室）或气泡（气泡室）：
径迹\emph{本身}就是一条催化成核的线。(d) 其
孔隙预先提供气泡核，从而防止过热以及
未加晶核的烧瓶中剧烈的“暴沸”。
\end{solution}

\begin{solution}{exo:b3:phase-transitions:10}
(a) 只有一个被打断的键：$2J$，与位置无关。(b)
$k_{\text{B}}\ln N$ 种排布方式。(c) $\Delta F = 2J -
k_{\text{B}}T\ln N \to -\infty$：在任何正温度下，
畴壁侵入，有序性随之消亡——$T_{\text{c}}(1\text{D}) = 0$。(d)
长度为 $\ell$ 的畴壁代价为 $2J\ell$，与熵
$\sim k_{\text{B}}\ell\ln3$ 相抗衡：当 $T < 2J/k_{\text{B}}\ln3 \approx
1.8\,J/k_{\text{B}}$ 时，长畴壁受到抑制，有序性得以保持
——量级恰好与昂萨格的精确结果
$2.27\,J/k_{\text{B}}$ 相当。
\end{solution}

\begin{solution}{exo:b3:phase-transitions:11}
(a) 两个\omterm{def:b3:phase-transitions:order}{序参量}都是选取符号的一个实数
（上/下；液体/气体关于临界密度）：属于同一对称性类。(b) 关联长度发散：
每个自旋的“平均场”与其自身一样涨落剧烈，
用邻居的平均值来代替已不再站得住脚。
(c) 近邻越多，平均越好；在四维以上，
涨落的作用范围增长得慢于平均化的改善，
\omterm{thm:b3:phase-transitions:meanfield}{平均场理论}变为精确。(d) 重整化群
反复平均掉最短的尺度，
把同一对称性、同一维数下的每一种微观模型
都流向同一个\omterm{def:b3:hamiltonian-mechanics:portrait}{不动点}；\omterm{prop:b3:phase-transitions:landau}{临界指数}是\omterm{def:b3:hamiltonian-mechanics:portrait}{不动点}的性质——
这正是水蒸气与磁铁
的小数位相同的原因。
\end{solution}

\begin{solution}{exo:b3:phase-transitions:12}
(a) 取 $\phi_0^2 = \alpha(T_{\text{c}} - T)/2b$，则 $F''(\phi_0) =
4\alpha(T_{\text{c}} - T) > 0$——沿径向的振荡
具有一个劲度系数，即（用场论的语言）一个质量，
且随势阱加深而增大。(b) 沿圆周的边缘处
势是平坦的：这是一个劲度系数为零、无质量的模式——
即戈德斯通模式，在磁体中表现为自旋波，在
超流体中表现为声子。(c) 真空中的 $\phi_0 \neq 0$ 即为凝聚体；
粒子质量即其与 $\phi_0$ 的耦合；径向振荡
即希格斯玻色子，已在 \qty{125}{GeV} 处找到。(d) 这种\omterm{def:b3:phase-transitions:order}{对称性破缺}、
凝聚体以及有质量与无质量的模式，最初都是实验室里的常见现象：
粒子物理引入的是一个固体物理的思想。
\end{solution}

\begin{solution}{pb:b3:phase-transitions:1}
\textbf{1.} $\mu_{\text{液体}}(T, P) = \mu_{\text{蒸气}}(T,
P)$：沸腾即共存，逐压强而论。
\textbf{2.} $\qty{3.6}{kPa/K}$。
\textbf{3.} $\ln2 = 4900(1/373 - 1/T)$：$T \approx \qty{394}{K} =
\qty{121}{\celsius}$。
\textbf{4.} 升高二十一度，按“每十度翻一倍”计，约
快四倍——这正是高压锅的全部商业逻辑。
\textbf{5.} $\ln0.65$：$T \approx \qty{361}{K} =
\qty{88}{\celsius}$——茶中单宁在低十二度时
萃取效果很差。
\textbf{6.} 溶质降低的是\emph{液体}的 $\mu$
（见\cref{exo:b3:grand-canonical:8}）；要与蒸气相匹配，
则需要更高的温度：沸点升高。
\textbf{7.} $\dd T/\dd P = \qty{-7.4e-8}{K/Pa}$：为负——
压强\emph{使冰熔化}。
\textbf{8.} 冰会漂浮：固相是密度较小的相，这种
反常（$\Delta v < 0$）使相界线向后倾斜——这种情形
很少见，因为大多数固体的堆积比其熔体更紧密。
\textbf{9.} $P = \rho gh \approx \qty{2.7}{MPa}$：仅使熔点变化
$-0.2$ K——但与地热相结合，它能维持冰床底部的
融水，成为冰川快速流动的润滑剂。
\textbf{10.} 摩擦生热与预熔的表面薄膜；
压强致熔的说法之所以长期流传，是因为斜率的符号确为真实——
只是其数值过小。
\textbf{11.} 冰会下沉，湖泊将从底部向上冻结
成固体块；没有隔热的冰盖，没有液态的避难所，
也没有淡水生物越冬的可能。
\textbf{12.} 十余种晶体态的冰意味着一幅\omterm{prop:b3:canonical-ensemble:machine}{自由能}
图景中有许多几乎简并的堆积方式：一个弯折的分子，
多种折衷——多形性是常态。
\textbf{13.} 两者都连续地缩小至零：在端点处，
两相合并，一级\omterm{def:b3:phase-transitions:order}{相变}线以
一个连续点的形式消失。
\textbf{14.} 在 $P > P_{\text{c}}$ 下加热至 $T > T_{\text{c}}$，
再减压：液体变为气体，全程无弯月面——
绕过山峰，而非翻越它。
\textbf{15.} 压缩系数发散，因此 $\sqrt N$ 量级的密度
涨落增长至光学尺度，强烈地散射光：
流体本身变成了一团云。
\textbf{16.} 超临界 CO$_2$ 从咖啡豆中萃取咖啡因；
泄压使其降到 $T_{\text{c}}$--$P_{\text{c}}$ 以下，而
“溶剂”便完全蒸发殆尽。
\textbf{17.} 在 \qty{221}{bar} 以上不存在液–气界面：
没有气泡，没有膜态沸腾，也没有烧毁危机——
热量流入一种光滑的单一流体。
\textbf{18.} \qty{400}{\celsius} 超过 $T_{\text{c}} =
\qty{374}{\celsius}$：在 \qty{300}{bar} 下，排出的水
处于超临界状态——沸腾并未被抑制，而是
\emph{无定义}。
\textbf{19.} 三条线交于三相点
（\qty{0.01}{\celsius}，\qty{611}{Pa}）；汽化线
向上延伸至其终点（\qty{647}{K}，\qty{221}{bar}）；
熔化线向左倾斜；升华发生在三相点以下。
\textbf{20.} 冷冻干燥在 \qty{611}{Pa} 以下进行，此时冰
升华；火星极冠因同样原因而季节性地升华。
\textbf{21.} 三相共存时\omterm{def:b3:lagrangian-mechanics:coordinates}{自由度}为零（吉布斯）：
状态为唯一的一点——这是测温学梦寐以求的，
却不可复现，因此长期用来定义开尔文。
\textbf{22.} 沿一条线，\omterm{def:b3:phase-transitions:order}{序参量}（密度）发生跃变，并
有潜热流动；在端点处，跃变已缩为零，
由涨落取而代之。
\textbf{23.} 若冷却表面将局部蒸气的
$\mu$ 拉得高于液体，水必然凝结：露、雾与
起雾的窗户，皆为 $\mu$ 交叉的可见表现。
\textbf{24.} 蒸发（液体 $\to$ 蒸气）发生在表面，
凝结（蒸气 $\to$ 液体）发生在锅盖上：上升携带的潜热
并未散入厨房，而是回送入锅中——
锅盖闭合了热循环。
\textbf{25.} $\dd P/\dd T = L/T\Delta v$：两个大气压下为 \qty{121}{\celsius}，
拉巴斯为 \qty{88}{\celsius}，熔化线
向后倾斜，而沸腾本身则终止于
（\qty{647}{K}，\qty{221}{bar}）——厨房，被绘制成了图。
\end{solution}
